\sectioncentered*{Заключение}
\addcontentsline{toc}{section}{Заключение}

В лабораторной работе разработано приложение с операциями над данными в оперативной памяти, выполнен анализ безопасности кода, использования ОЗУ и процессорного времени, проанализированы дампы после создания, обновления и удаления записей, построена методика оценки и проведено сравнение с предоставленным вариантом.

\textbf{Семантическая составляющая.} Безопасность кода определяется не только применением криптографии, но и временем жизни открытого текста в ОЗУ, наличием затирания буферов и явным разделением конфиденциальных и неконфиденциальных сущностей. Дампы памяти делают эти свойства наблюдаемыми: заметки как открытый текст могут оставаться в дампе; в модели секрета после обработки лежит только шифртекст, при этом кратковременные копии строк CPython в дампе ещё возможны, в том числе после удаления (критерий C7).

\textbf{Прагматическая составляющая.} Получены воспроизводимый стенд для CRUD в ОЗУ, сценарий снятия дампов и чек-лист методики, позволяющий сравнивать реализации и формулировать улучшения: сокращение копий строк, использование защищённых буферов, контроль журналирования и прав доступа к ключу шифрования.

\begin{thebibliography}{9}
	\addcontentsline{toc}{section}{Список использованных источников}

	\bibitem{howard}
	Ховард М., Лебланк Д. Защищённый код. \mdash  М. : Русская редакция, 2004. \mdash  704~с.

	\bibitem{howard24}
	Ховард М., Лебланк Д., Вьега Дж. 24 смертных греха компьютерной безопасности. \mdash  СПб. : Питер, 2010. \mdash  400~с.

	\bibitem{cprofile}
	The Python Profilers (\texttt{cProfile}) [Электронный ресурс]. \mdash  Режим доступа: \url{https://docs.python.org/3/library/profile.html}. \mdash  Дата доступа: 12.09.2026.

	\bibitem{gcore}
	gcore(1) \mdash  Generate a core file of a running program [Электронный ресурс]. \mdash  Режим доступа: \url{https://man7.org/linux/man-pages/man1/gcore.1.html}. \mdash  Дата доступа: 10.09.2026.

	\bibitem{fernet}
	Fernet (symmetric encryption) [Электронный ресурс]. \mdash  Режим доступа: \url{https://cryptography.io/en/latest/fernet/}. \mdash  Дата доступа: 10.09.2026.
\end{thebibliography}
